\documentclass[12pt]{article}
\usepackage[utf8]{inputenc}
\usepackage[T1]{fontenc}
\usepackage{hyperref}
\title{Tafi Language Resources for K-12 Curriculum Development}
\author{Compiled by Claude Sonnet 5 (claude-sonnet-5) for Andr\'as Kornai}
\date{}
\begin{document}
\maketitle
\section{Introduction}
Tafi (ISO 639-3 tcd) is a Ghana-Togo Mountain language spoken in four communities of the Hohoe District, Volta Region, Ghana. Ewe is the dominant lingua franca and has influenced Tafi. Search results cite about 4,400 speakers (2013), far fewer than the figure in the task description. Do not confuse it with the closely related Nyangbo (Tutrugbu).
\section{Primary Educational Resources}
No primers, textbooks or school programmes were found in search results.
\section{Linguistic and Reference Materials}
Mercy Bobuafor's 2013 Leiden doctoral thesis \emph{A Grammar of Tafi} is the first comprehensive description (phonology, morphology, syntax, glossed folktales and proverbs):
\url{https://www.universiteitleiden.nl/en/research/research-output/humanities/a-grammar-of-tafi}

Leiden repository copy: \url{https://scholarlypublications.universiteitleiden.nl/access/item%3A2922605/view}

Overview of the related Nyangbo-Tafi language: \url{https://en.wikipedia.org/wiki/Nyangbo-Tafi_language}

A project abstract mentions a lexicon and multimedia corpus, but whether a Tafi dictionary was published is not confirmed in the results.
\section{Religious and Community Resources}
None found in search results.
\section{Digital and Audio Resources}
No standalone audio or software resources found. The grammar project mentions a multimedia corpus; its availability was not verified.
\section{Curriculum Development Context}
Tafi is not known to be taught in schools. The grammar could underpin an orthography and teaching materials, but mother-tongue materials would have to be created, probably alongside Ewe and English.
\end{document}